% !TeX root = ../../book.tex
\section{强归纳法}
\secbegin{secStrongInduction}
\index{induction!strong|(}

考虑下面的例子，我们将尝试用归纳法来证明它。

\begin{example}
\label{exProofByWeakInductionFail}
递归地定义一个序列
\[ b_0 = 1 \quad \text{and} \quad b_{n+1} = 1 + \sum_{k=0}^n b_k \text{ 对于所有 } n \in \mathbb{N} \]
我们将尝试通过归纳法证明对于所有 $n \in \mathbb{N}$, $b_n = 2^n$。

\begin{itemize}
\item (\textbf{基本情况}) 根据序列的定义可得 $b_0=1=2^0$。目前为止，一切正常。

\item (\textbf{递归步骤}) 给定 $n \in \mathbb{N}$，并假设 $b_n = 2^n$。我们需要证明 $b_{n+1}=2^{n+1}$。

而 $b_{n+1} = 1 + \sum_{k=0}^n b_k = {\dots}$ 悲剧了。
\end{itemize}

问题出在这里。如果我们可以将求和中的每个 $b_k$ 替换为 $2^k$，那么我们就能完成证明。然而我们不能证明这种替换是合理的：我们的归纳假设只给我们提供了关于 $b_n$ 的信息，而不是关于 $k < n$ 的通用表示 $b_k$。
\end{example}

\textit{强}归纳原理与弱归纳原理看起来很像，只是它的归纳假设更强大。尽管名字里带强，但强归纳法并不比弱归纳更强。这两个原理是等效的。事实上，我们将\textit{通过弱归纳法}来证明强归纳法原理！

\begin{theorem}[强归纳法原理]
\label{thmStrongInduction}
\index{induction!on $\mathbb{N}$ (strong)}
\index{strong induction principle}
设 $p(n)$ 为逻辑公式其中 $n \in \mathbb{N}$ 为自由变量，并设  $n_0 \in \mathbb{N}$。如果
\begin{enumerate}[(i)]
\item $p(n_0)$ 为真；且
\item 对于所有 $n \ge n_0$，如果 $p(k)$ 对于所有 $n_0 \le k \le n$ 为真，则 $p(n+1)$ 为真；
\end{enumerate}
那么对于所有 $n \ge n_0$, $p(n)$ 为真。
\end{theorem}

\begin{cproof}
对于每个 $n \ge n_0$，设 $q(n)$ 为 $p(k)$ 对于所有 $n_0 \le k \le n$ 为真的断言。

%% BEGIN EXTRACT (xtrAssumingImplicationsExample) %%
注意 $q(n)$ 意味着对于所有 $n \ge n_0$, $p(n)$，因为给定 $n \ge n_0$，如果 $p(k)$ 对于所有 $n_0 \le k \le n$ 为真，则当 $k=n$ 时 $p(k)$ 为真。

因此只要证明 $q(n)$ 对于所有 $n \ge n_0$ 都为真即可。
%% END EXTRACT %%
我们通过弱归纳法来证明这一点。

\begin{itemize}
\item (\textbf{基本情况}) $q(n_0)$ 等价于 $p(n_0)$，因为唯一具有 $n_0 \le k \le n_0$ 的自然数 $k$ 就是 $n_0$ 本身；因此 $q(n_0)$ 根据条件 (i) 为真。

\item (\textbf{归纳步骤}) 设 $n \ge n_0$ 并假设 $q(n)$ 为真。则 $p(k)$ 对于所有 $n_0 \le k \le n$ 都为真。

我们需要证明 $q(n+1)$ 为真 --- 即 $p(k)$ 对于所有 $n_0 \le k \le n+1$ 为真。我们已知 $p(k)$ 对于所有 $n_0 \le k \le n$ 为真 --- 此为归纳假设 --- 则根据条件 (ii) $p(n+1)$ 为真。因此我们有 $p(k)$ 对于所有 $n_0 \le k \le n+1$ 都为真。
\end{itemize}
根据归纳法，$q(n)$ 对于所有 $n \ge n_0$ 为真。因此，$p(n)$ 对于所有 $n \ge n_0$ 都为真。
\end{cproof}

\begin{strategy}[用强归纳法证明]
\label{strStrongInduction}
\index{proof!by strong induction}
为了证明 $\forall n \ge n_0,\, p(n)$ 形式的命题，只需证明：
\begin{itemize}
\item (\textbf{基本情况}) $p(n_0)$ 为真；
\item (\textbf{归纳步骤}) 对于所有 $n \ge n_0$，如果 $p(k)$ 对于所有 $n_0 \le k \le n$ 为真，则 $p(n+1)$ 为真。
\end{itemize}
\end{strategy}

与弱归纳法一样，我们可以用图表说明强归纳法的工作原理。归纳假设用大方块表示，包含所有 $n_0 \le k \le n$ 下的 $p(k)$，其中 $p(n_0)$ 是基本情况。

\begin{center}
\begin{tikzpicture}
\draw[fill=indstepcol] (-1,-1) rectangle (8,1);
\draw (0,0) node[diamond, draw, text width=27, text centered, fill=indbasecol] {$n_0$};
\draw (2,0) node[circle, draw, text width=27, text centered, fill=white] {$n_0+1$};
\draw (3.5,0) node {$\cdots$};
\draw (5,0) node[circle, draw, text width=27, text centered, fill=white] {$n-1$};
\draw (7,0) node[circle, draw, text width=27, text centered, fill=white] {$n$};
\draw (9.5,0) node[circle, draw, dashed, text width=27, text centered, fill=white](nplusone) {$n+1$};
\draw[->] (8,0) -- (nplusone);
\end{tikzpicture}
\end{center}

可见与弱归纳法的唯一区别在于归纳假设。
\begin{itemize}
\item \textbf{弱归纳步骤：} 给定 $n \ge n_0$，\fbox{假设 $p(n)$ 为真}~，推导 $p(n+1)$；
\item \textbf{强归纳步骤：} 给定 $n \ge n_0$，\fbox{假设 $p(k)$ 对于所有 $n_0 \le k \le n$ 为真}~，推导 $p(n+1)$。
\end{itemize}

现在我们用强归纳法来完成\Cref{exProofByWeakInductionFail}的证明。

\begin{example}[\Cref{exProofByWeakInductionFail} revisited]
递归地定义一个序列
\[ b_0 = 1 \quad \text{and} \quad b_{n+1} = 1 + \sum_{k=0}^n b_k \text{ for all } n \in \mathbb{N} \]
我们将通过强归纳法证明，对于所有 $n \in \mathbb{N}$, $b_n = 2^n$。

\begin{itemize}
\item (\textbf{基本情况}) 根据序列的定义可得 $b_0=1=2^0$。
\item (\textbf{归纳步骤}) 给定 $n \in \mathbb{N}$，并假设对于所有 $k \le n$, $b_k=2^k$。我们需要证明 $b_{n+1}=2^{n+1}$。这是成立的，因为
\begin{align*}
b_{n+1} &= 1 + \sum_{k=0}^n b_k && \text{根据 $b_{n+1}$ 的递归公式} \\
&= 1 + \sum_{k=0}^n 2^k && \text{根据递归假设} \\
&= 1 + (2^{n+1}-1) && \text{根据 \Cref{exSumOfPowersOf2}} \\
&= 2^{n+1}
\end{align*}
\end{itemize}
通过归纳法可得对于所有 $n \in \mathbb{N}$, $b_n=2^n$。
\end{example}

以下定理将强归纳原理应用于证明，其中我们需要在归纳步骤中参考\textit{固定}数量的先前步骤。

\begin{theorem}[强归纳原理 {(多基本情况)}]
\label{thmSIPMultipleBaseCases}
设 $p(n)$ 为逻辑公式，其中 $n \in \mathbb{N}$ 为自由变量，并设 $n_0 < n_1 \in \mathbb{N}$。如果
\begin{enumerate}[(i)] \vspace{5pt} 
\item $p(n_0), p(n_0+1), \dots, p(n_1)$ 都为真；且
\item 对于所有 $n \ge n_1$，如果 $p(k)$ 对于 $n_0 \le k \le n$ 为真，则 $p(n+1)$ 为真；
\end{enumerate}
\vspace{5pt}
那么 $p(n)$ 对于所有 $n \ge n_0$ 为真。
\end{theorem}

\begin{cproof}
$p(n)$ 对于所有 $n \ge n_0$ 都为真可以通过强归纳法得出的：
\begin{itemize}
\item (\textbf{基本情况}) 由 (i) 可得 $p(n_0)$ 为真；
\item (\textbf{归纳步骤}) 给定 $n \ge n_0$ 并假设 $p(k)$ 对于所有 $n_0 \le k \le n$为真。则：
\begin{itemize}
\item 如果 $n < n_1$，则 $n+1 \le n_1$，所以根据 (i) $p(n)$ 为真；
\item 如果 $n \ge n_1$，则根据 (ii) $p(n+1)$ 为真。
\end{itemize}
上述两种情况我们都得到 $p(n+1)$ 为真。
\end{itemize}
因此，根据强归纳法，我们有 $p(n)$ 对于所有 $n \ge n_0$ 都为真。
\end{cproof}

\begin{strategy}[通过具有多基本情况的强归纳法证明]
\index{proof!by strong induction with multiple base cases}
要证明 $\forall n \ge n_0,\, p(n)$ 形式的陈述，只需证明：
\begin{itemize}
\item (\textbf{基本情况}) 对于所有 $k \in \{ n_0, n_0+1, \dots, n_1 \}$, $p(k)$，其中 $n_1 > n_0$；且
\item (\textbf{归纳步骤}) 对于所有 $n \ge n_1$，如果 $p(k)$ 对于所有 $n_0 \le k \le n$ 为真，那么 $p(n+1)$ 为真。
\end{itemize}
\end{strategy}

这种强归纳法与通常的强归纳法有两处不同：
\begin{itemize}
\item 存在多个基本情况 $p(n_0), p(n_0+1), \dots, p(n_1)$，不是只有一个。
\item 归纳步骤即涉及最小情况 $n_0$ 又涉及最大情况 $n_1$：归纳步骤中的变量 $n$ 被设为大于等于 $n_1$，而归纳假设假设对于所有 $n_0 \le k \le n$, $p(k)$。
\end{itemize}

下图展示了多基本情况的强归纳是如何工作的。

\begin{center}
\begin{tikzpicture}
\draw[fill=indstepcol] (-1,-1) rectangle (9,1);
\draw (0,0) node[diamond, draw, text width=27, text centered, fill=indbasecol] {$n_0$};
\draw (1.5,0) node {$\cdots$};
\draw (3,0) node[diamond, draw, text width=27, text centered, fill=indbasecol] {$n_1$};
\draw (5,0) node[circle, draw, text width=27, text centered, fill=white] {$n_1+1$};
\draw (6.5,0) node {$\cdots$};
\draw (8,0) node[circle, draw, text width=27, text centered, fill=white] {$n$};
\draw (10.5,0) node[circle, draw, dashed, text width=27, text centered, fill=white](nplusone) {$n+1$};
\draw[->] (9,0) -- (nplusone);
\end{tikzpicture}
\end{center}

请注意常规强归纳法和具有多基本情况的强归纳法之间在归纳步骤中变量量化的差异：
\begin{itemize}
\item \textbf{单基本情况 } 给定 $n \ge \boxed{n_0}$ 并假设 $p(k)$ 对于所有 $\boxed{n_0} \le k \le n$ 为真。
\item \textbf{多基本情况 } 给定 $n \ge \boxed{n_1}$ 并假设 $p(k)$ 对于所有 $\boxed{n_0} \le k \le n$ 为真。
\end{itemize}

在强归纳步骤中变量 $n$ 和 $k$ 的量化至关重要，因为量化会影响对 $n$ 和 $k$ 的假设。

当证明有关递归定义的序列的结果时，通常会出现对多基本情况的需求，其中通项的定义取决于固定数量的前项的值。

\begin{example}
定义序列
\[ a_0 = 0, \quad a_1 = 1, \quad a_n = 3a_{n-1} - 2a_{n-2} \text{ for all } n \ge 2 \]
找到并证明用 $n$ 表示 $a_n$ 的通项公式。下面给出了序列的前几项，建立了一个我们可以尝试证明的模式：
\begin{center}
\begin{tabular}{c|ccccccccc}
$n$   & $0$ & $1$ & $2$ & $3$ & $4$  & $5$  & $6$  & $7$   & $8$ \\ \hline
$a_n$ & $0$ & $1$ & $3$ & $7$ & $15$ & $31$ & $63$ & $127$ & $255$
\end{tabular}
\end{center}

对于所有 $n \ge 0$ 通项公式似乎是 $a_n = 2^n - 1$。我们将 $n=0$ 和 $n=1$ 的情况作为基本情况，通过强归纳法证明这一点。

\begin{itemize}
\item (\textbf{基本情况}) 根据序列的定义，可得：
\begin{itemize}
\item $a_0 = 0 = 2^0 - 1$;
\item $a_1 = 1 = 2^1 - 1$;
\end{itemize}
因此，当 $n=0$ 和 $n=1$ 时，该断言为真。

\item (\textbf{归纳步骤}) 给定 $n \ge 1$ 并假设对于所有 $0 \le k \le n$, $a_k = 2^k - 1$。我们需要证明 $a_{n+1} = 2^{n+1} - 1$。

因为 $n \ge 1$，可得 $n+1 \ge 2$，所以我们可以将递归公式应用于 $a_{n+1}$。即
\begin{align*}
a_{n+1} &= 3a_n - 2a_{n-1} && \text{根据 $a_{n+1}$ 的定义} \\
&= 3(2^n-1) - 2(2^{n-1}-1) && \text{根据递归假设} \\
&= 3 \cdot 2^n - 3 - 2 \cdot 2^{n-1} + 2 && \text{展开} \\
&= 3 \cdot 2^n - 3 - 2^n + 2 && \text{指数定律} \\
&= 2 \cdot 2^n - 1 && \text{化简} \\
&= 2^{n+1} - 1 && \text{指数定律}
\end{align*}
\end{itemize}
最后由归纳法得出结论。
\end{example}

以下练习通过多基本情况的强归纳法进行证明。

\begin{exercise}
对于所有 $n \ge 2$，递归地定义序列 $a_0 = 4$, $a_1 = 9$ 和 $a_n = 5a_{n-1} - 6a_{n-2}$。证明对于所有 $n \in \mathbb{N}$, $a_n = 3 \cdot 2^n + 3^n$。
\end{exercise}

\begin{exercise}
\textit{Tribonacci 序列}为序列 $t_0, t_1, t_2, \dots$ 定义为
\[ t_0 = 0, \quad t_1 = 0, \quad t_2 = 1, \quad t_n = t_{n-1} + t_{n-2} + t_{n-3} \text{ for all } n \ge 3 \]
证明对于所有 $n \ge 3$, $t_n \le 2^{n-3}$。
\end{exercise}

\begin{exercise}
\textit{Frobenius硬币问题}问的是何时可以从给定面值的硬币中凑出给定数量的货币。例如，仅使用 $3$ 面值和 $5$ 面值的硬币无法获得 $7$ 面值，但对于所有 $n \ge 8$，仅使用 $3$ 面值和 $5$ 面值的硬币是\textit{可以}获得 $n$ 面值的。请证明这一点。
\hintlabel{exFrobeniusCoinProblem}{%
请注意，如果仅使用 $3$ 面值和 $5$ 面值的硬币即可获得 $n$ 面值，那么 $n+3$ 也可以获得。想一下如何利用这一点引到多基本情况的强归纳法。
}
\end{exercise}

\subsection*{良序原则}

我们可以对自然数进行归纳和递归的根本原因是它们的排序方式。具体来说，自然数满足\textit{良序原则}：如果一个自然数集合至少有一个元素，则它有一个最小元素。这将自然数与其他数集区分开来；例如，$\mathbb{Z}$ 就没有最小元素，因为如果 $a \in \mathbb{Z}$ 则 $a-1 \in \mathbb{Z}$ 且 $a-1 < a$。

\begin{theorem}[良序原则]
\label{thmWellOrderingPrinciple}
\index{well-ordering principle}
设 $X$ 为自然数集合。如果 $X$ 非空，则 $X$ 有一个最小元素。
\end{theorem}

\begin{cidea}
假设 $X$ 是一个自然数集合，我们要证明的命题具有 $p \Rightarrow q$ 的形式。即
\[ X \text{ 非空} \quad \Rightarrow \quad X \text{ 有一个最小元素} \]
假设 $X$ 非空并没有给我们太多帮助，所以让我们尝试一下反证：
\[ X \text{ 没有最小元素} \quad \Rightarrow \quad X \text{ 为空} \]
$X$ 没有最小元素的假设确实给了我们一些可用的东西。在这个假设下，我们需要推导出 $X$ 是空的。

我们将利用强归纳法`逼迫 $X$ 增大'来做到这一点。当然 $0 \not \in X$，否则它就是最小元素。如果数字 $0, 1, \dots, n$ 都不是 $X$ 的元素，那么 $n+1$ 也不能是，因为如果是，那么\textit{它}将是 $X$ 的最小元素。让我们把这个论点正式化。
\end{cidea}

\begin{cproof}
设 $X$ 为一个没有最小元素的自然数集合。我们用强归纳法证明对于所有 $n \in \mathbb{N}$, $n \not \in X$。
\begin{itemize}
\item (\textbf{基本情况}) $0 \not \in X$ 因为如果 $0 \in X$ 则 $0$ 必然是 $X$ 的最小元素，因为 0 是最小的自然数。
\item (\textbf{归纳步骤}) 假设对于所有 $0 \le k \le n$, $k \not \in X$。如果 $n+1 \in X$ 则 $n+1$ 为 $X$ 的最小元素；实际上，如果 $\ell < n+1$ 则 $0 \le \ell \le n$，所以根据递归假设 $\ell \not \in X$。这与 $X$ 没有最小元素的假设相矛盾，因此 $n+1 \not \in X$。
\end{itemize}
根据强归纳法，对于每个 $n \in \mathbb{N}$, $n \not \in X$。因为 $X$ 是一个自然数集合，并且它不包含自然数，因此 $X$ 为空集。
\end{cproof}

下面证明 $\sqrt{2}$ 是无理数的证明是良序原理的经典应用。

\begin{proposition}
\label{propSqrt2Irrational}
$\sqrt{2}$ 为无理数。
\end{proposition}

要证明 \Cref{propSqrt2Irrational} 我们需要使用如下引理，其利用良序原理证明分数可以`消去最低项'。

\begin{lemma}
\label{lemFractionInLowestTerms}
设 $q$ 为正有理数。存在一对非零自然数 $a,b$ 满足 $q=\frac{a}{b}$ 并且能够整除 $a$ 和 $b$ 的唯一自然数是 $1$。
\end{lemma}

\begin{cproof}
首先请注意，我们可以将 $q$ 表示为两个非零自然数的比，因为 $q$ 是一个正有理数。根据良序原则，存在一个\textit{最小}自然数 $a$，使得对于某正 $b \in \mathbb{N}$, $q=\frac{a}{b}$。

假设 $d$ 为除 $1$ 之外的某自然数，可以整除 $a$ 和 $b$。请注意 $d \ne 0$，因为如果 $d=0$ 则意味着 $a=0$。由于$d \ne 1$，因此 $d \ge 2$。

因为 $d$ 整除 $a$ 和 $b$，则存在自然数 $a',b'$ 满足 $a=a'd$ 且 $b=b'd$。此外，$a',b'>0$ 因为 $a,b,d > 0$。可得
\[ q = \frac{a}{b} = \frac{a'd}{b'd} = \frac{a'}{b'} \]
而 $d \ge 2$，因此
\[ a' = \frac{a}{d} \le \frac{a}{2} < a \]
这与 $a$ 的最小性相矛盾。因此，我们假设除 $1$ 之外还存在某个自然数 $d$ 可以整除 $a$ 和 $b$ 是错误的 --- 因此，唯一可以整除 $a$ 和 $b$ 的自然数是 $1$。
\end{cproof}

现在我们可以开始证明 $\sqrt{2}$ 为无理数了。

\begin{cproof}[of \Cref{propSqrt2Irrational}]
假设 $\sqrt{2}$ 为有理数。因为 $\sqrt{2}>0$，这意味着可以写成
\[ \sqrt{2} = \frac{a}{b} \]
其中 $a$ 和 $b$ 都是正自然数。根据\Cref{lemFractionInLowestTerms}，我们可以假设唯一能整除 $a$ 和 $b$ 的自然数为 $1$。

将方程两边平方并变成乘法得
\[ a^2=2b^2 \]
因为 $a^2$ 为偶数。根据\Cref{propIfProductEvenThenSomeFactorEven}，$a$ 为偶数，所以可以写成 $a=2c$ 其中 $c > 0$。 那么 $a^2=(2c)^2=4c^2$，因此
\[ 4c^2=2b^2 \]
两边除以 $2$ 得
\[ 2c^2 = b^2 \]
因此 $b^2$ 为偶数。根据\Cref{propIfProductEvenThenSomeFactorEven}, $b$ 为偶数。

然而如果 $a$ 和 $b$ 都为偶数，自然数 $2$ 可以整除 $a$ 和 $b$。这与只有唯一自然数 $1$ 可以整除 $a$ 和 $b$ 矛盾。因此假设中 $\sqrt{2}$ 为有理数错误，$\sqrt{2}$ 为无理数。
\end{cproof}

\begin{writingtip}
在 \Cref{propSqrt2Irrational} 的证明中，我们可以分别证明 $a^2$ 为偶数意味着 $a$ 是偶数，并且 $b^2$ 为偶数意味着 $b$ 是偶数。这需要我们重复两次相同的证明，有点繁琐且乏味！单独证明辅助结论（如\Cref{lemFractionInLowestTerms}），然后在后面的定理中引用它们可以提高核心证明的可读性，还有助于强调重要步骤，特别是当辅助结论特别技术性时。
\end{writingtip}

\begin{exercise}
如果我们尝试证明 $\sqrt{4}$ 是无理数，那么 \Cref{propSqrt2Irrational} 的证明会出现什么问题？
\end{exercise}

\begin{exercise}
\label{exSquareRootThreeIsIrrational}
证明 $\sqrt{3}$ 为无理数。
\begin{backhint}
\hintref{exSquareRootThreeIsIrrational}
首先证明如果 $a \in \mathbb{Z}$ 且 $a^2$ 能被$3$ 整除，那么 $a$ 能被 $3$ 整除。
\end{backhint}
\end{exercise}


\index{induction!strong|)}